\begin{abstract}
Can a language model turn a short description into a three-dimensional assembly? In \textsc{Text-to-Brick}, a generator must infer dimensions and proportions, place discrete pieces, and account for stud connectivity. We present BrickForma, which transfers geometric knowledge from large language models through executable parametric programs. These programs generate 80,000 voxel models in 21 categories on a $48^3$ grid, with five description levels per shape. Same-color tiling, connection analysis, and local repair produce checked brick targets for fine-tuning Qwen3-1.7B. At inference, a mixed-integer program seeks an exact cover of the predicted voxel union with one stud-connected component. On 265 matched prompts, BrickForma's raw outputs attain the highest mean VQAScore (0.1079), PE (0.1572), and SigLIP2 (0.0565) among the three systems. BrickGPT leads VQAScore at the shortest description level and has a lower mean invalid-placement ratio (0.00\% versus 0.81\%). On 110 earlier raw/retiled pairs, the number of single-component layouts rises from 18 to 110, and all pairs have equal occupied voxel unions after translation.
\end{abstract}
